% Options for packages loaded elsewhere
\PassOptionsToPackage{unicode}{hyperref}
\PassOptionsToPackage{hyphens}{url}
\documentclass[
  ukrainian,
]{article}
\usepackage{xcolor}
\usepackage[margin=2.5cm]{geometry}
\usepackage{amsmath,amssymb}
\setcounter{secnumdepth}{-\maxdimen} % remove section numbering
\usepackage{iftex}
\ifPDFTeX
  \usepackage[T1]{fontenc}
  \usepackage[utf8]{inputenc}
  \usepackage{textcomp} % provide euro and other symbols
\else % if luatex or xetex
  \usepackage{unicode-math} % this also loads fontspec
  \defaultfontfeatures{Scale=MatchLowercase}
  \defaultfontfeatures[\rmfamily]{Ligatures=TeX,Scale=1}
\fi
\usepackage{lmodern}
\ifPDFTeX\else
  % xetex/luatex font selection
\fi
% Use upquote if available, for straight quotes in verbatim environments
\IfFileExists{upquote.sty}{\usepackage{upquote}}{}
\IfFileExists{microtype.sty}{% use microtype if available
  \usepackage[]{microtype}
  \UseMicrotypeSet[protrusion]{basicmath} % disable protrusion for tt fonts
}{}
\makeatletter
\@ifundefined{KOMAClassName}{% if non-KOMA class
  \IfFileExists{parskip.sty}{%
    \usepackage{parskip}
  }{% else
    \setlength{\parindent}{0pt}
    \setlength{\parskip}{6pt plus 2pt minus 1pt}}
}{% if KOMA class
  \KOMAoptions{parskip=half}}
\makeatother
\usepackage{longtable,booktabs,array}
\newcounter{none} % for unnumbered tables
\usepackage{calc} % for calculating minipage widths
% Correct order of tables after \paragraph or \subparagraph
\usepackage{etoolbox}
\makeatletter
\patchcmd\longtable{\par}{\if@noskipsec\mbox{}\fi\par}{}{}
\makeatother
% Allow footnotes in longtable head/foot
\IfFileExists{footnotehyper.sty}{\usepackage{footnotehyper}}{\usepackage{footnote}}
\makesavenoteenv{longtable}
\usepackage{graphicx}
\graphicspath{{/Users/kalinicenkovitalijmikolajovic/.codex/worktrees/5c15/Script_for_translate_book/Output_Final/farlou_chunked_v1_pdf_chunks/}}
\makeatletter
\newsavebox\pandoc@box
\newcommand*\pandocbounded[1]{% scales image to fit in text height/width
  \sbox\pandoc@box{#1}%
  \Gscale@div\@tempa{\textheight}{\dimexpr\ht\pandoc@box+\dp\pandoc@box\relax}%
  \Gscale@div\@tempb{\linewidth}{\wd\pandoc@box}%
  \ifdim\@tempb\p@<\@tempa\p@\let\@tempa\@tempb\fi% select the smaller of both
  \ifdim\@tempa\p@<\p@\scalebox{\@tempa}{\usebox\pandoc@box}%
  \else\usebox{\pandoc@box}%
  \fi%
}
% Set default figure placement to htbp
\def\fps@figure{htbp}
\makeatother
\ifLuaTeX
\usepackage[bidi=basic,shorthands=off,provide=*]{babel}
\else
\usepackage[bidi=default,shorthands=off,provide=*]{babel}
\fi
\ifLuaTeX
  \usepackage{selnolig} % disable illegal ligatures
\fi
\setlength{\emergencystretch}{3em} % prevent overfull lines
\providecommand{\tightlist}{%
  \setlength{\itemsep}{0pt}\setlength{\parskip}{0pt}}
\usepackage{bookmark}
\IfFileExists{xurl.sty}{\usepackage{xurl}}{} % add URL line breaks if available
\urlstyle{same}
\hypersetup{
  pdflang={uk},
  hidelinks,
  pdfcreator={LaTeX via pandoc}}

\author{}
\date{}

\begin{document}

\section{Как играть в «Блуждающего
пьяницу»?}\label{ux43aux430ux43a-ux438ux433ux440ux430ux442ux44c-ux432-ux431ux43bux443ux436ux434ux430ux44eux449ux435ux433ux43e-ux43fux44cux44fux43dux438ux446ux443}

ШАГ 1. Блуждания пьяницы начинаются из произвольной точки сетки (в нашем
случае это точка A).

\pandocbounded{\includegraphics[keepaspectratio]{images/_page_321_Figure_4.jpeg}}

\begin{itemize}
\tightlist
\item
  А начальная точка - - внутренние точки сетки - □ -Конечная точка р;
\item
  \texttt{g\_i} вознаграждение в m . \texttt{p\_I}
\end{itemize}

РИС. 43.1. Игровое поле для игры «Блуждающий пьяница».

ШАГ 2. На каждом ходе игры пьяница случайным образом перебирается в одну
из четырех соседних точек сетки (в нашем случае соседними для A являются
точки B, C, D и E и вероятность попадания в каждую из этих точек равна
0,25).

Таблица 43.1

Вероятность окончания случайного блуждания в граничной точке
\texttt{p\_i} и величина вознаграждения \texttt{g\_i}

{\def\LTcaptype{none} % do not increment counter
\begin{longtable}[]{@{}
  >{\raggedright\arraybackslash}p{(\linewidth - 8\tabcolsep) * \real{0.2523}}
  >{\raggedright\arraybackslash}p{(\linewidth - 8\tabcolsep) * \real{0.4541}}
  >{\raggedright\arraybackslash}p{(\linewidth - 8\tabcolsep) * \real{0.2752}}
  >{\raggedright\arraybackslash}p{(\linewidth - 8\tabcolsep) * \real{0.0092}}
  >{\raggedright\arraybackslash}p{(\linewidth - 8\tabcolsep) * \real{0.0092}}@{}}
\toprule\noalign{}
\begin{minipage}[b]{\linewidth}\raggedright
Граничная точка р і
\end{minipage} & \begin{minipage}[b]{\linewidth}\raggedright
\texttt{P\_\{\textbackslash{}mathcal\{A\}\}\textbackslash{}left(p\_\{i\}\textbackslash{}right)}
--- относительная доля блужданий, закончившихся в точке
\texttt{p\_\{i\}}
\end{minipage} & \begin{minipage}[b]{\linewidth}\raggedright
\texttt{g\_i} --- величина вознаграждения при достижении точки
\texttt{p\_i}
\end{minipage} & \begin{minipage}[b]{\linewidth}\raggedright
\end{minipage} & \begin{minipage}[b]{\linewidth}\raggedright
\end{minipage} \\
\midrule\noalign{}
\endhead
\bottomrule\noalign{}
\endlastfoot
1234567891011 & 0,040,150,030,060,170,050,060,150,030,060,160,04 &
110000000 & & \\
\end{longtable}
}

ti

ШАГ 3. После перехода в соседнюю точку процесс блуждания возобновляется.
Так пьяница идет от точки к точке до тех пор, пока случайным образом не
окажется в граничной точке \texttt{p\_t} . Здесь он останавливается, а
мы записываем номер точки \texttt{p\_t} . На этом заканчивается одна
случайная прогулка.

ШАГ 4. Повторим щаги 1-3 достаточно много раз и найдем относительное
число прибытий пьяницы в каждую из граничных точек \texttt{p\_i} . В
табл. 43.1 приведены типичные результаты, полученные в результате 100
случайных прогулок.

\texttt{ШАГ\ Б.\ Предположим,\ что\ пьяница\ получает\ вознаграждение\ \ \textquotesingle{}g\_i\ =\ g\textbackslash{}left(p\_i\textbackslash{}right)\textquotesingle{}\ \ (величина\ граничного\ условия\ в\ точке\ \ \textquotesingle{}p\_i\textquotesingle{}\ ,\ если\ он\ достиг\ точки\ \ \textquotesingle{}p\_i\textquotesingle{}\ )\ и\ предположим,\ что\ цель\ игры—вычислить\ среднее\ вознаграждение\ \ \textquotesingle{}R\textbackslash{}left(A\textbackslash{}right)\textquotesingle{}\ \ для\ всех\ случайных\ прогулок,\ начинающихся\ из\ точки\ A.\ Этот\ средний\ выигрыш\ определяется\ по\ формуле}

\begin{verbatim}
Игра завершается, когда величина  R\left(A\right)  найдена. В соответствии с данными табл. 43.1 для  R\left(A\right)  получаем
\end{verbatim}

Теперь объясним, зачем нам понадобилось играть в такую игру.

\section{Для чего играть в «Блуждающего
пьяницу»?}\label{ux434ux43bux44f-ux447ux435ux433ux43e-ux438ux433ux440ux430ux442ux44c-ux432-ux431ux43bux443ux436ux434ux430ux44eux449ux435ux433ux43e-ux43fux44cux44fux43dux438ux446ux443}

Оказывается, что среднее вознаграждение, о котором только что шла речь,
является приближенным решением задачи Дирихле в точке A. Это интересное
наблюдение основано на двух фактах.

\begin{enumerate}
\def\labelenumi{\arabic{enumi}.}
\item
  Предположим, что пъяница начал свою прогулку из точки, лежащей на
  еранице. Каждая такая прогулка немедленно заканчивается в той же
  точке, и пъяница получает вознаграждение \texttt{g\_i} . Таким
  образом, среднее вознаграждение для каждой граничной точки равно
  \texttt{g\_i} .
\item
  Теперь предположим, что прогулка нач нается из внутренней точки. Тогда
  ясно, что среднее вознаграждечие для точки \texttt{R(A)} будет средним
  арифметическим от средних вознаграждений для четырех соседних точек
\end{enumerate}

\begin{verbatim}
\begin{cases} u(x, 0) = \sin(\pi x), & 0 \le x \le 1, \\ u_t(x, 0) = 0. & 0 \le x \le 1, \end{cases}
\end{verbatim}

Зададимся еще раз вопросом: почему же среднее значение вознаграждения
\texttt{R(A)} аппроксимирует решение задачи Дирихле в точке A? Мы
видели, что величина \texttt{R(A)} удовлетворяет двум

уравнениям

(во внутренних точках), \texttt{R(A)} = \texttt{g\_i} (в граничных
точках).

Если \texttt{g\_i} ---это значения функции \texttt{g(x,\ y)} из
граничного условия в граничных точках \texttt{p\_i} , то два наших
уравнения точно совпадают с двумя уравнениями, которые были получены при
решении задачи Дирихле методом конечных разностей. То есть величина
\texttt{R(A)} соответствует величине \texttt{u\_\{i,f\}} в разностных
уравнениях

(i, j)---внутренняя точка, \texttt{u\_\{ij\}\ =\ g\_\{ij\},}
\texttt{g\_\{ij\}} ---значение решения в граничной точке (i, j).

Следовательно, величина \texttt{R(A)} действительно аппроксимирует
решение уравнения с частными производными в точке A.

Теперь можно считать, что игра в «Блуждающего пьяницу» состоит из трех
шагов.

\section{Решение уравнения Лапласа методом
Монте-Карло}\label{ux440ux435ux448ux435ux43dux438ux435-ux443ux440ux430ux432ux43dux435ux43dux438ux44f-ux43bux430ux43fux43bux430ux441ux430-ux43cux435ux442ux43eux434ux43eux43c-ux43cux43eux43dux442ux435-ux43aux430ux440ux43bux43e}

\begin{itemize}
\tightlist
\item
  Приведенные ниже правила позволяют получить решение в одной внутренней
  точке квадрата.
\item
  ШАГ 1. Совершим некоторое число случайных прогулок, начинающихся в
  точке А и заканчивающихся в одной из граничных точек. Проследим за
  тем, сколько раз прогулки заканчивались в каждой граничной точке.
\item
  ШАГ 2. После того как все прогулки завершены, для каждой граничной
  точки вычисляем относительное число прогулок, завершившихся в этой
  точке. Обозначим эти величины \texttt{P\_A(p\_i)} .

  \begin{itemize}
  \tightlist
  \item
    ШАГ 3. Вычисляем приближенное решение \texttt{u(A)} по формуле
    \texttt{u(A)\ =\ g\_1\ P\_A(p\_1)\ +\ g\_2\ P\_A(p\_2)\ +\ \textbackslash{}dots\ +\ g\_N\ P\_A(p\_N)}
    ,
  \end{itemize}
\end{itemize}

где \texttt{g\_i} --- значение функции g в точке \texttt{p\_i} , а N ---
число граничных точек.

Игру «Блуждающий пьяница» можно модернизировать, чтобы решить более
сложную задачу. Ниже приводится пример такой модернизации.

Решение задачи Дирихле с переменными коэффициентами.

Рассмотрим эллиптическую краевую вадачу в квадрате

Дайте физическую интерпретацию этой задачи.

\begin{enumerate}
\def\labelenumi{\arabic{enumi}.}
\setcounter{enumi}{3}
\tightlist
\item
  Для уравнения \texttt{u\_\{tt\}\ =\ u\_\{xx\}}найдите все решения вида
\end{enumerate}

\texttt{u\_\{xx\}\ +\ (\textbackslash{}sin\ x)\ u\_\{yy\}\ =\ 0} ,
\texttt{0\ \textless{}\ x\ \textless{}\ \textbackslash{}pi} ,
\texttt{0\ \textless{}\ y\ \textless{}\ \textbackslash{}pi} ,
\texttt{(\textbackslash{}Gamma\ Y)} \texttt{u(x,\ y)\ =\ g(x,\ y)} на
границе квадрата.

Чтобы решить эту задачу, заменим \texttt{u\_\{xx\}} , \texttt{u\_\{yy\}}
и \texttt{\textbackslash{}sin\ x} следующим образом:

\$\$

Будет ли решением сумма двух решений такого вида?

\end{document}
